% GENERATED FILE. Edit canonical lessons, then run npm run generate:books.
% canonical-source-hash: fnv1a64:c5d8aa44e3a3b017
% canonical-lessons: KA-C393-samsattu, KA-C393-mantri, KA-C393-samvidhana, KA-C393-prajaprabhutva, KA-C393-pratinidhisu

\chapter{Running a Country}
\label{ch:ka-running-a-country}
\hlchaptermodality{393}

\begin{chapteropening}
\textbf{By the end of this chapter:} \emph{I can say a parliament, a minister, a constitution, democracy and to represent.}

Complete the last lesson of chapter 393: I can say a parliament, a minister, a constitution, democracy and to represent.
\end{chapteropening}

\section[\texorpdfstring{\kn{ಸಂಸತ್ತು} (saṃsattu)}{ಸಂಸತ್ತು (saṃsattu)}]{\kn{ಸಂಸತ್ತು} (saṃsattu) — a parliament}

\label{lesson:KA-C393-samsattu}



\begin{quote}
Before the new one: say the Kannada for to campaign, then the Kannada for to address a crowd.
\end{quote}

\subsection*{You'll want to know: \kn{ಸಂಸತ್ತು}}
\textbf{\kn{ಸಂಸತ್ತು}} — \emph{saṃsattu} — \textquotedblleft{}a parliament\textquotedblright{}.

\begin{quote}
\kn{ಹೊಸ} \kn{ನಿಯಮದ} \kn{ಬಗ್ಗೆ} \kn{ಸಂಸತ್ತಿನಲ್ಲಿ} \kn{ಚರ್ಚೆ} \kn{ನಡೆಯಿತು}.
\end{quote}

\begin{quote}
\emph{hosa niyamada bagge saṃsattinalli carce naḍeyitu.}
\end{quote}

\begin{quote}
There was a debate in parliament on the new rule.
\end{quote}

\kn{ಸಂಸದ}, a member of parliament.

\begin{cousinweb}
\textbf{\kn{ಸಂಸತ್ತು}} (\emph{saṃsattu}) is borrowed from Sanskrit \emph{saṃsad}, \textquotedblleft{}a sitting together, an assembly\textquotedblright{} (Monier-Williams), with the Kannada -\emph{u} after the final consonant.
\end{cousinweb}

\begin{grammarlens}[title={What you've built}]
One more word for giving reasons and opinions.
\end{grammarlens}

\subsection*{Your turn}
\begin{itemize}
\raggedright
  \item \emph{Say it:} \emph{saṃsattu}
  \item \emph{Say it:} \emph{saṃsattu}, once more
  \item \emph{Say it:} \emph{hosa niyamada bagge saṃsattinalli carce naḍeyitu.} — There was a debate in parliament on the new rule.
  \item \emph{Recall:} say the Kannada for to campaign, then the Kannada for to address a crowd, then say \emph{saṃsattu} again
\end{itemize}

\begin{tcolorbox}[breakable,colback=teal!4,colframe=teal!35!black,title={Before you move on}]
What does \kn{ಸಂಸತ್ತು} mean? (\textquotedblleft{}A parliament\textquotedblright{}.) Say it once more.
\end{tcolorbox}
\section[\texorpdfstring{\kn{ಮಂತ್ರಿ} (mantri)}{ಮಂತ್ರಿ (mantri)}]{\kn{ಮಂತ್ರಿ} (mantri) — a minister}

\label{lesson:KA-C393-mantri}



\begin{quote}
Before the new one: say the Kannada for to address a crowd, then the Kannada for a parliament.
\end{quote}

\subsection*{You'll want to know: \kn{ಮಂತ್ರಿ}}
\textbf{\kn{ಮಂತ್ರಿ}} — \emph{mantri} — \textquotedblleft{}a minister\textquotedblright{}.

\begin{quote}
\kn{ಮಂತ್ರಿ} \kn{ನಮ್ಮ} \kn{ಶಾಲೆಗೆ} \kn{ಭೇಟಿ} \kn{ಕೊಟ್ಟರು}.
\end{quote}

\begin{quote}
\emph{mantri namma śālege bhēṭi koṭṭaru.}
\end{quote}

\begin{quote}
The minister visited our school.
\end{quote}

Plural: \kn{ಮಂತ್ರಿಗಳು}, \emph{mantrigaḷu}.

Respectful plural verb: \kn{ಮಂತ್ರಿ} \kn{ಬಂದರು}.

\begin{cousinweb}
\textbf{\kn{ಮಂತ್ರಿ}} (\emph{mantri}) is borrowed from Sanskrit \emph{mantrin}, \textquotedblleft{}one who gives counsel\textquotedblright{}, a king's adviser (CDIAL 9838).
\end{cousinweb}

\begin{grammarlens}[title={What you've built}]
One more word for giving reasons and opinions.
\end{grammarlens}

\subsection*{Your turn}
\begin{itemize}
\raggedright
  \item \emph{Say it:} \emph{mantri}
  \item \emph{Say it:} \emph{mantri}, once more
  \item \emph{Say it:} \emph{mantri namma śālege bhēṭi koṭṭaru.} — The minister visited our school.
  \item \emph{Recall:} say the Kannada for to address a crowd, then the Kannada for a parliament, then say \emph{mantri} again
\end{itemize}

\begin{tcolorbox}[breakable,colback=teal!4,colframe=teal!35!black,title={Before you move on}]
What does \kn{ಮಂತ್ರಿ} mean? (\textquotedblleft{}A minister\textquotedblright{}.) Say it once more.
\end{tcolorbox}
\section[\texorpdfstring{\kn{ಸಂವಿಧಾನ} (saṃvidhāna)}{ಸಂವಿಧಾನ (saṃvidhāna)}]{\kn{ಸಂವಿಧಾನ} (saṃvidhāna) — a constitution}

\label{lesson:KA-C393-samvidhana}



\begin{quote}
Before the new one: say the Kannada for a parliament, then the Kannada for a minister.
\end{quote}

\subsection*{You'll want to know: \kn{ಸಂವಿಧಾನ}}
\textbf{\kn{ಸಂವಿಧಾನ}} — \emph{saṃvidhāna} — \textquotedblleft{}a constitution\textquotedblright{}.

\begin{quote}
\kn{ಸಂವಿಧಾನ} \kn{ಎಲ್ಲರಿಗೂ} \kn{ಸಮಾನ} \kn{ಹಕ್ಕು} \kn{ಕೊಡುತ್ತದೆ}.
\end{quote}

\begin{quote}
\emph{saṃvidhāna ellarigū samāna hakku koḍuttade.}
\end{quote}

\begin{quote}
The constitution gives everyone equal rights.
\end{quote}

\kn{ಸಂವಿಧಾನಬದ್ಧ}, constitutional.

\begin{cousinweb}
\textbf{\kn{ಸಂವಿಧಾನ}} (\emph{saṃvidhāna}) is borrowed from Sanskrit \emph{saṃvidhāna}, \textquotedblleft{}arrangement, management\textquotedblright{} (Monier-Williams).
\end{cousinweb}

\begin{grammarlens}[title={What you've built}]
One more word for giving reasons and opinions.
\end{grammarlens}

\subsection*{Your turn}
\begin{itemize}
\raggedright
  \item \emph{Say it:} \emph{saṃvidhāna}
  \item \emph{Say it:} \emph{saṃvidhāna}, once more
  \item \emph{Say it:} \emph{saṃvidhāna ellarigū samāna hakku koḍuttade.} — The constitution gives everyone equal rights.
  \item \emph{Recall:} say the Kannada for a parliament, then the Kannada for a minister, then say \emph{saṃvidhāna} again
\end{itemize}

\begin{tcolorbox}[breakable,colback=teal!4,colframe=teal!35!black,title={Before you move on}]
What does \kn{ಸಂವಿಧಾನ} mean? (\textquotedblleft{}A constitution\textquotedblright{}.) Say it once more.
\end{tcolorbox}
\section[\texorpdfstring{\kn{ಪ್ರಜಾಪ್ರಭುತ್ವ} (prajāprabhutva)}{ಪ್ರಜಾಪ್ರಭುತ್ವ (prajāprabhutva)}]{\kn{ಪ್ರಜಾಪ್ರಭುತ್ವ} (prajāprabhutva) — democracy}

\label{lesson:KA-C393-prajaprabhutva}



\begin{quote}
Before the new one: say the Kannada for a minister, then the Kannada for a constitution.
\end{quote}

\subsection*{You'll want to know: \kn{ಪ್ರಜಾಪ್ರಭುತ್ವ}}
\textbf{\kn{ಪ್ರಜಾಪ್ರಭುತ್ವ}} — \emph{prajāprabhutva} — \textquotedblleft{}democracy\textquotedblright{}.

\begin{quote}
\kn{ಪ್ರಜಾಪ್ರಭುತ್ವದಲ್ಲಿ} \kn{ಜನರೇ} \kn{ಸರ್ಕಾರ} \kn{ಆರಿಸುತ್ತಾರೆ}.
\end{quote}

\begin{quote}
\emph{prajāprabhutvadalli janarē sarkāra ārisuttāre.}
\end{quote}

\begin{quote}
In a democracy the people choose the government.
\end{quote}

\kn{ಪ್ರಜೆ} (citizen) + \kn{ಪ್ರಭುತ್ವ} (rule).

\begin{cousinweb}
\textbf{\kn{ಪ್ರಜಾಪ್ರಭುತ್ವ}} (\emph{prajāprabhutva}) is a modern word built from Sanskrit \emph{prajā}, \textquotedblleft{}people, subjects\textquotedblright{}, and \emph{prabhutva}, \textquotedblleft{}lordship, rule\textquotedblright{} (Monier-Williams): rule by the people.
\end{cousinweb}

\begin{grammarlens}[title={What you've built}]
One more word for giving reasons and opinions.
\end{grammarlens}

\subsection*{Your turn}
\begin{itemize}
\raggedright
  \item \emph{Say it:} \emph{prajāprabhutva}
  \item \emph{Say it:} \emph{prajāprabhutva}, once more
  \item \emph{Say it:} \emph{prajāprabhutvadalli janarē sarkāra ārisuttāre.} — In a democracy the people choose the government.
  \item \emph{Recall:} say the Kannada for a minister, then the Kannada for a constitution, then say \emph{prajāprabhutva} again
\end{itemize}

\begin{tcolorbox}[breakable,colback=teal!4,colframe=teal!35!black,title={Before you move on}]
What does \kn{ಪ್ರಜಾಪ್ರಭುತ್ವ} mean? (\textquotedblleft{}Democracy\textquotedblright{}.) Say it once more.
\end{tcolorbox}
\section[\texorpdfstring{\kn{ಪ್ರತಿನಿಧಿಸು} (pratinidhisu)}{ಪ್ರತಿನಿಧಿಸು (pratinidhisu)}]{\kn{ಪ್ರತಿನಿಧಿಸು} (pratinidhisu) — to represent}

\label{lesson:KA-C393-pratinidhisu}



\begin{quote}
Before the new one: say the Kannada for a constitution, then the Kannada for democracy.
\end{quote}

\subsection*{You'll want to know: \kn{ಪ್ರತಿನಿಧಿಸು}}
\textbf{\kn{ಪ್ರತಿನಿಧಿಸು}} — \emph{pratinidhisu} — \textquotedblleft{}to represent\textquotedblright{}.

\begin{quote}
\kn{ಅವಳು} \kn{ರಾಜ್ಯವನ್ನು} \kn{ಸ್ಪರ್ಧೆಯಲ್ಲಿ} \kn{ಪ್ರತಿನಿಧಿಸಿದಳು}.
\end{quote}

\begin{quote}
\emph{avaḷu rājyavannu spardheyalli pratinidhisidaḷu.}
\end{quote}

\begin{quote}
She represented the state in the competition.
\end{quote}

Past (I …): \kn{ಪ್ರತಿನಿಧಿಸಿದೆ}, \emph{pratinidhiside}.

\kn{ಪ್ರತಿನಿಧಿ}, a representative.

\begin{cousinweb}
\textbf{\kn{ಪ್ರತಿನಿಧಿಸು}} (\emph{pratinidhisu}) is built on Sanskrit \emph{pratinidhi}, \textquotedblleft{}a substitute, a representative\textquotedblright{} (Monier-Williams), with the Kannada verb ending -\emph{isu}.
\end{cousinweb}

\begin{grammarlens}[title={What you've built}]
One more word for giving reasons and opinions.
\end{grammarlens}

\subsection*{Your turn}
\begin{itemize}
\raggedright
  \item \emph{Say it:} \emph{pratinidhisu}
  \item \emph{Say it:} \emph{pratinidhisu}, once more
  \item \emph{Say it:} \emph{avaḷu rājyavannu spardheyalli pratinidhisidaḷu.} — She represented the state in the competition.
  \item \emph{Recall:} say the Kannada for a constitution, then the Kannada for democracy, then say \emph{pratinidhisu} again
\end{itemize}

\begin{tcolorbox}[breakable,colback=teal!4,colframe=teal!35!black,title={Before you move on}]
What does \kn{ಪ್ರತಿನಿಧಿಸು} mean? (\textquotedblleft{}To represent\textquotedblright{}.) Say it once more.
\end{tcolorbox}
